\documentclass{article}

\usepackage{amsmath,amssymb}
\usepackage{xurl}
\usepackage[hidelinks]{hyperref}
\setlength{\emergencystretch}{2em}

% Shared commands for notes in docs/paper-gaps/.

\DeclareUnicodeCharacter{2080}{\ifmmode{}_{0}\else\textsubscript{0}\fi}
\DeclareUnicodeCharacter{2081}{\ifmmode{}_{1}\else\textsubscript{1}\fi}
\DeclareUnicodeCharacter{2082}{\ifmmode{}_{2}\else\textsubscript{2}\fi}
\DeclareUnicodeCharacter{2099}{\ifmmode{}_{n}\else\textsubscript{n}\fi}

\newcommand{\Lean}{\textsc{Lean}}
\ExplSyntaxOn
\NewDocumentCommand{\leanid}{m}{
  \ifmmode
    \textsf{\detokenize{#1}}
  \else
    \group_begin:
    \urlstyle{sf}
    \tl_set:Nn \l_tmpa_tl {#1}
    \tl_replace_all:Nnn \l_tmpa_tl {\_} {_}
    \exp_args:NV \path \l_tmpa_tl
    \group_end:
  \fi
}
\ExplSyntaxOff
\providecommand{\tr}{\operatorname{tr}}
\newcommand{\tnleanIssueBase}{https://github.com/LionSR/TNLean/issues/}
\newcommand{\tnleanPrBase}{https://github.com/LionSR/TNLean/pull/}
\newcommand{\tnleanDocBase}{https://sirui-lu.com/TNLean/docs/}
\newcommand{\ghissue}[1]{\href{\tnleanIssueBase#1}{GitHub issue \##1}}
\newcommand{\ghpr}[1]{\href{\tnleanPrBase#1}{PR \##1}}
\newcommand{\doclink}[1]{\href{\tnleanDocBase#1.html}{\path{#1}}}

% Dirac notation and the identity operator, as in blueprint/src/macros/common.tex.
% \providecommand keeps note-local definitions authoritative.
\usepackage{dsfont}
\providecommand{\ket}[1]{|#1\rangle}
\providecommand{\bra}[1]{\langle #1|}
\providecommand{\braket}[2]{\langle #1 | #2 \rangle}
\providecommand{\ketbra}[2]{|#1\rangle\!\langle #2|}
\providecommand{\Id}{\mathds{1}}
\providecommand{\one}{\mathds{1}}
\providecommand{\C}{\mathbb{C}}
\providecommand{\MN}[1]{M_{#1}(\C)}
\providecommand{\MD}{M_D(\C)}

% External referencing for paper sources.  The build helper writes sanitized
% auxiliary files under build/paper-aux/ and excludes duplicated source labels.
\usepackage{xr}

% Import a generated paper auxiliary file with a caller-supplied label prefix.
% The candidate paths support builds from docs/paper-gaps/, the repository root,
% and build/paper-aux/ itself.
\newcommand{\importpaperaux}[2]{%
  \IfFileExists{../../build/paper-aux/#2.aux}{%
    \externaldocument[#1]{../../build/paper-aux/#2}%
  }{%
    \IfFileExists{build/paper-aux/#2.aux}{%
      \externaldocument[#1]{build/paper-aux/#2}%
    }{%
      \IfFileExists{#2.aux}{%
        \externaldocument[#1]{#2}%
      }{}%
    }%
  }%
}

% General external reference macros
\newcommand{\paperref}[2]{\ref{#1-#2}}
\newcommand{\papereqref}[2]{(\ref{#1-#2})}

% CPSV16 source (1606.00608 / MPDO-22-12-17-2.tex) external document and shortcuts
\importpaperaux{cpsv16-}{1606.00608/MPDO-22-12-17-2-tnlean}
\newcommand{\cpsvref}[1]{\paperref{cpsv16}{#1}}
\newcommand{\cpsveqref}[1]{\papereqref{cpsv16}{#1}}

% Verdict marker: \gapnote{<kind>}{<status>}. Kind is the policy
% classification (clarification, local-correction, scope-restriction,
% unfaithful, false-source, open-gap); status is open, wip (elimination
% underway), resolved, or historical. Machine-read by the site index;
% prints a verdict line.
\newcommand{\gapnote}[2]{\par\noindent\textbf{Verdict:} #1 (#2).\par}


\title{The Historical Range Restriction in the $Y_1$--$X_2$ Mixed Contraction}
\author{}
\date{2026-08-10}

\begin{document}
\maketitle
\gapnote{clarification}{open}

\section{The graphical step}

A historical formalization attempt replaced the paper gate
$u=Y_2\mathbin{-}Y_1$ by the mixed contraction $Y_1\mathbin{-}X_2$ and then
contracted a normalized open tail between source-cut factors. This replacement
does not occur in Ref.~\cite[lines~536--557]{Cirac2017MPU}. Interpreted
as an ambient matrix identity, this simplification would identify the
open-tail coefficient with the identity on the full product-index space. That
identity is false in general because the source factors need not be
surjective. The available isometry relation is only
\begin{align}
    X_2^\dagger X_2
        &= 1.
    \label{eq:mpu_mixed_kernel_range_restriction_source_isometry}
\end{align}
Equation~\ref{eq:mpu_mixed_kernel_range_restriction_source_isometry} does not imply
the ambient coisometry relation $X_2X_2^\dagger=1$.

\section{Range-restricted interpretation}

Let $M_1$ and $M_2$ denote the source-cut factors, let $X_1$ and $X_2$ denote
the corresponding source isometries, and define
$W_\rho=\rho\otimes 1_d$. The valid interpretation retains the external
$Y_1$--$X_2$ mixed kernels and closes the complete tensor network before
moving the cut. The two projection identities used in this closed network are
\begin{align}
    X_1X_1^\dagger W_\rho M_1
        &= M_1,
    \label{eq:mpu_mixed_kernel_range_restriction_first_range_identity}\\
    X_2X_2^\dagger M_2
        &= M_2.
    \label{eq:mpu_mixed_kernel_range_restriction_second_range_identity}
\end{align}
These identities follow from the source-cut factorizations and the recovery
formulas for $Y_1$ and $Y_2$. They assert that the relevant source-cut factors
lie in the ranges of $X_1$ and $X_2$; neither identity asserts that a source
isometry is an ambient coisometry.

The forward internal kernel also depends on orientation. It consists of an
output-first double-layer letter followed by the $K$th power of its normalized
diagonal. Conjugate transposition reverses the virtual chain. Consequently, the
reflected simple-contraction equations cannot be substituted pointwise into
this coefficient.

The historical route would next compare two complete contractions: the
ordinary $Y_1$--$X_2$ mixed-kernel Gram matrix and the normalized closed output
tail on $K+2$ sites. The normalized open-tail metric does not establish this
equality, and the comparison is not an obligation for the paper gate $u$.

\section{Verdict}

The range identities are valid auxiliary algebra, but the surrounding
$Y_1$--$X_2$ network is not the source gate of Ref.~\cite{Cirac2017MPU}.
Accordingly, this note no longer supplies a route to the paper's source-$u$
isometry lemma. The paper-facing proof must start from $u=Y_2\mathbin{-}Y_1$.

The forward kernel and its output-layer expression are formalized by
\leanid{MPOTensor.sourceY₁X₂ForwardKernel} and
\leanid{MPOTensor.sourceY₁X₂ForwardKernel_eq_outputLayer_mul_normalizedDiagonal_pow}.
The range identities in
Eqs.~\ref{eq:mpu_mixed_kernel_range_restriction_first_range_identity} and
\ref{eq:mpu_mixed_kernel_range_restriction_second_range_identity} are formalized by
\leanid{MPOTensor.sourceX₁_mul_conjTranspose_mul_weight_mul_sourceCutM₁} and
\leanid{MPOTensor.sourceX₂_mul_conjTranspose_mul_sourceCutM₂}, respectively.
The fully expanded normalized metric is formalized by
\leanid{MPOTensor.IsMPU.normalized_sourceY₁X₂_openTail_metric}. This theorem does
not identify the normalized metric with the ordinary $Y_1$--$X_2$ mixed-kernel
Gram matrix, and no such equality is claimed as part of the paper-gate proof.

\bibliographystyle{alpha}
\bibliography{references}

\end{document}
